\capitulo{7}{Conclusiones y trabajo futuro}

\section{Conclusiones}

En el presente proyecto se ha diseñado, implementado y evaluado de forma preliminar un sistema de monitorización y análisis para Blender. La solución desarrollada se plantea como una infraestructura de software modular orientada a la captura automática de datos técnicos durante sesiones de modelado tridimensional, su almacenamiento estructurado en formato CSV y su posterior análisis mediante métricas y visualizaciones integradas en el propio entorno de Blender.

Durante la fase de investigación y desarrollo se ha constatado la viabilidad técnica de registrar determinados cambios y estados internos de Blender sin requerir una intervención continua por parte del usuario. La solución opera en segundo plano durante la sesión de modelado y permite recopilar información sobre eventos, transformaciones, cambios de modo, variaciones geométricas y otros indicadores asociados al proceso de trabajo. No obstante, esta captura debe entenderse dentro de las limitaciones propias de la API de Blender, ya que no todos los eventos internos son accesibles de forma directa ni homogénea. Por este motivo, el sistema combina distintas estrategias, como la comparación de estados, el uso de temporizadores, la detección de operadores recientes y el análisis posterior de los registros generados.

El sistema implementado ha permitido estructurar información temporal, espacial y operativa en un formato legible y reutilizable. A partir de los datos registrados, el add-on de análisis calcula métricas que ayudan a interpretar aspectos del proceso de modelado, como la evolución temporal de la sesión, los cambios en la escena, las modificaciones sobre los objetos y determinadas propiedades geométricas o UV. Estos resultados no deben interpretarse como una medición completa o definitiva del comportamiento del usuario, sino como indicadores técnicos que pueden apoyar el análisis del proceso seguido durante una sesión de modelado.

Un aspecto relevante del trabajo ha sido la incorporación de métricas geométricas y UV, incluyendo el análisis de características de malla, la evaluación de islas UV y la comparación entre modelos mediante medidas como la distancia de Hausdorff. Estas funcionalidades amplían el alcance del sistema más allá del simple registro de acciones, permitiendo obtener información sobre el estado y la evolución de los modelos tridimensionales. Sin embargo, su interpretación requiere cautela, ya que la calidad técnica de un modelo o el nivel de experiencia de un usuario no pueden deducirse únicamente a partir de estas métricas. Para utilizar estos indicadores con fines evaluativos sería necesario complementar el sistema con criterios expertos, mayor validación empírica y comparación con datos de referencia.

La arquitectura modular adoptada, basada en dos complementos independientes coordinados mediante un formato CSV común, ha resultado adecuada para separar la fase de adquisición de datos de la fase de análisis posterior. Esta separación reduce el acoplamiento entre componentes, facilita la revisión independiente de los registros y permite reutilizar los datos capturados en distintos procesos de análisis. Como parte de la validación funcional, se ha utilizado un CSV real procedente de la actividad realizada en la EASD, Escuela de Arte y Superior de Diseño de Burgos, con el fin de comprobar el flujo completo de captura, carga, análisis y visualización. Este registro debe entenderse como evidencia funcional sobre datos reales, no como una evaluación experimental completa con usuarios. Al mismo tiempo, el uso de CSV introduce ciertas limitaciones, como la necesidad de mantener la coherencia del esquema, controlar la compatibilidad entre versiones y validar correctamente los datos antes de calcular métricas.

La interfaz desarrollada para el add-on de análisis permite presentar métricas complejas de forma más organizada mediante paneles, agrupaciones visuales y ayudas contextuales. Esta integración dentro de Blender facilita que el usuario pueda revisar los resultados sin abandonar el entorno de trabajo. Aun así, la experiencia de uso debe considerarse mejorable, especialmente en lo relativo a la claridad del inicio de la captura, la validación de que una sesión se está registrando correctamente y la comunicación de errores o advertencias durante el proceso.

En conjunto, el proyecto demuestra que es posible construir una herramienta funcional para capturar, almacenar, analizar y visualizar datos asociados al proceso de modelado 3D en Blender. Los resultados obtenidos permiten considerar la solución como un prototipo académico viable y extensible, con potencial para ser utilizado en contextos educativos, experimentales o de investigación sobre flujos de trabajo en modelado tridimensional. No obstante, no debe entenderse todavía como un producto final ni como una herramienta plenamente validada para escenarios profesionales o de producción.

Para alcanzar un mayor grado de madurez sería necesario ampliar la validación con más usuarios y sesiones, medir cuantitativamente el impacto en el rendimiento, evaluar la precisión de la detección de eventos, reforzar la separación entre la lógica independiente y la capa dependiente de Blender, mejorar la reproducibilidad de las pruebas automatizadas y profundizar en los aspectos de privacidad y tratamiento de datos. Estas líneas de mejora permitirían consolidar la herramienta y aumentar la confianza en los resultados obtenidos.

\section{Conclusiones técnicas}

A nivel de ingeniería de software, el desarrollo del proyecto ha aportado evidencias funcionales preliminares respecto al comportamiento de las API de gráficos en tiempo real y la gestión de la carga computacional en hilos de interfaz. Estas observaciones proceden de la implementación y de pruebas controladas, pero no sustituyen a una campaña experimental amplia con usuarios finales:

\begin{itemize}
    \item \textbf{Aprovechamiento del Grafo de Dependencias}: Se ha comprobado de forma funcional que la suscripción de funciones analíticas a los disparadores pos-evaluación del grafo (\texttt{depsgraph\_update\_post}) permite detectar modificaciones estructurales relevantes de la escena, requiriendo un control estricto de cortocircuitos lógicos para discriminar eventos redundantes.
    \item \textbf{Mitigación del impacto algorítmico}: Uno de los mayores desafíos técnicos consistió en preservar la tasa de refresco gráfica (\textit{framerate}) del visor 3D. El diseño utiliza firmas de estado compactas y limitación de la frecuencia de muestreo mediante temporizadores a intervalos de $0.25$ segundos. La medición cuantitativa del impacto en CPU y fluidez queda como una validación futura necesaria para confirmar esta mejora en distintos equipos y escenas.
    \item \textbf{Desacoplamiento del renderizado analítico}: Se implementó el aislamiento de los cálculos estadísticos respecto a la capa gráfica. Mediante el uso del motor abstracto \texttt{FigureCanvasAgg} de \textit{matplotlib}, el sistema procesa los diagramas en un búfer de memoria volátil e inyecta los resultados como matrices de píxeles puros en texturas procedurales. Este enfoque bajo demanda reduce el acoplamiento, aunque su impacto en tiempos de respuesta debe medirse de forma sistemática.
    \item \textbf{Gestión de materiales dinámicos}: La instanciación en tiempo de ejecución de mallas poligonales mediante estructuras \texttt{bmesh} y la configuración automática de materiales translúcidos basados en nodos de emisión y mezclas (\textit{Mix Shader}) se han empleado como soluciones funcionales para representar analíticas visuales superpuestas en el espacio de trabajo. Su legibilidad debe evaluarse con usuarios y escenas diversas antes de considerarse una solución definitiva.
\end{itemize}

\section{Limitaciones del sistema}

A pesar de que el sistema desarrollado cumple rigurosamente con los requisitos funcionales planteados en los objetivos, se han identificado las siguientes limitaciones técnicas que delimitan su alcance actual:

\begin{itemize}
    \item \textbf{Dependencia de la plataforma anfitriona}: Debido a su estrecho acoplamiento con la API de Python de Blender (\texttt{bpy}), el núcleo del sistema no es directamente portable a otras suites de creación tridimensional o herramientas CAD de la industria (como Autodesk Maya o 3ds Max) sin una reingeniería completa de la capa de adquisición.
    \item \textbf{Saturación de almacenamiento en sesiones extendidas}: Aunque la captura diferencial reduce la redundancia, las sesiones de modelado de larga duración pueden generar un volumen considerable de datos tabulares, lo que exige políticas de archivado, compactación o segmentación de archivos CSV.
    \item \textbf{Sensibilidad a la escala jerárquica}: El cálculo de distancias de similitud geométrica mediante la distancia de Hausdorff muestra una alta dependencia de la correcta aplicación de las transformaciones de escala del objeto. La presencia de transformaciones no aplicadas en la matriz mundial puede introducir desviaciones marginales en las métricas si el usuario trabaja fuera del origen de coordenadas.
\end{itemize}

\section{Líneas de trabajo futuras}

Partiendo de la base arquitectónica establecida por el sistema, se proponen las siguientes líneas de investigación y desarrollo para la expansión del sistema:

\begin{itemize}
    \item \textbf{Integración de modelos de aprendizaje automático}: Incorporar algoritmos de minería de datos y aprendizaje no supervisado (\textit{clustering}) para agrupar automáticamente los históricos de interacción, permitiendo clasificar de forma predictiva los perfiles de los usuarios (experto, intermedio o novato) en función de sus patrones secuenciales de modelado.
    \item \textbf{Optimización del motor de persistencia}: Sustituir la lectura secuencial en streaming de archivos CSV por un motor de base de datos relacional embebido y ligero (como SQLite) o sistemas de almacenamiento binario indexado, optimizando la velocidad de acceso y consulta en análisis comparativos masivos.
    \item \textbf{Abstracción de una capa de API universal}: Rediseñar el módulo de captura de datos mediante un patrón de diseño abstracto que independice la lógica analítica del software anfitrión, facilitando la creación de conectores homólogos para otras herramientas estándares del sector como Unity, Unreal Engine o Maya.
    \item \textbf{Internacionalización y personalización de la interfaz}: Ampliar el sistema de textos de interfaz para permitir alternar entre español e inglés, ajustar unidades y adaptar la presentación de métricas a distintos perfiles de usuario.
    \item \textbf{Sistemas de asistencia interactiva y tutorización inteligente}: Desarrollar un motor de inferencia en tiempo real que, a partir de las métricas de errores geométricos frecuentes (como vértices duplicados o normales invertidas), actúe como un tutor proactivo. Este módulo podrá predecir fallos topológicos inminentes y ofrecer sugerencias contextuales al estudiante, maximizando el valor pedagógico del sistema en entornos educativos técnicos.
\end{itemize}

En conclusión, el presente proyecto demuestra la viabilidad técnica de unificar la adquisición no intrusiva y la visualización analítica en el espacio tridimensional. El sistema se consolida como un prototipo funcional para registrar y analizar procesos de modelado en Blender, pero su uso como herramienta de auditoría del comportamiento de diseño o de evaluación educativa requiere una validación experimental posterior con usuarios, protocolos definidos y métricas de rendimiento reproducibles.
